\documentclass[11pt,letterpaper]{article}
\usepackage{rf}

\begin{document}

\title{Proofs}
\author{Leo Lee}
\date{\today}

\maketitle

\section{Introduction}

This section shows why we can approximate functions as Taylor polynomials while retaining some extent of stability.

To seek for $[SO(3) \times SO(2)]$-invariant Ricci solitons given by $(M, g)$ and function $M \to \R$, where $M$ is $\R^3 \times S^1$ or $\R^2 \times S^2$, we shall study the solutions of the following ODE system:
\begin{gather}
    f'' = \frac{a''}{a} + 2 \frac{b''}{b} - 1, \\
    a'' = -2\frac{a'b'}{b} + a'f' + a, \\
    b'' = \frac{1 - (b')^2}{b} - \frac{a'b'}{a} + b'f' + b.
\end{gather}
for functions
\[ a: [0, \infty) \to [0, \infty), \quad b: [0, \infty) \to [0, \infty), \quad f: [0, \infty) \to \R. \]

\section{Rigorous Taylor Remainder Bounds via Majorants}

\subsection{Power-Series Recurrence}

Let
\[
a(r)=\sum_{k=0}^{\infty}a_k r^k,\qquad
b(r)=\sum_{k=0}^{\infty}b_k r^k,\qquad
f(r)=\sum_{k=0}^{\infty}f_k r^k
\]
be formal power-series solutions of the given ODE system.
Substitution into the denominator-cleared equations and
collection of coefficients of equal powers of \(r\) yield
recurrence relations of the form
\[
A_k
\begin{pmatrix}
a_{k+2}\\
b_{k+2}\\
f_{k+2}
\end{pmatrix}
=
F_k(a_0,\ldots,a_{k+1},
    b_0,\ldots,b_{k+1},
    f_0,\ldots,f_{k+1}),
\]
where \(A_k\) is the coefficient matrix of the highest-order
unknowns and \(F_k\) contains the remaining terms.

The recurrence is well-defined whenever \(A_k\) is invertible.
Its precise form, including all contributions to \(A_k\), must
be obtained by coefficient extraction from the original ODEs.

\subsection{Majorant Theorem}

Suppose that the initial coefficients and the recurrence
relations satisfy the following conditions for constants
\(K>0\) and \(C>0\):

\begin{enumerate}
    \item The base coefficients satisfy
    \[
    |a_j|,\ |b_j|,\ |f_j|\le KC^j
    \]
    for every index \(j\) required to initialize the induction.

    \item For every \(k\ge 0\), whenever
    \[
    |a_j|,\ |b_j|,\ |f_j|\le KC^j,
    \qquad 0\le j\le k+1,
    \]
    the recurrence implies
    \[
    |a_{k+2}|,\ |b_{k+2}|,\ |f_{k+2}|
    \le KC^{k+2}.
    \]
\end{enumerate}

Then, by strong induction,
\[
\boxed{
|a_k|,\ |b_k|,\ |f_k|\le KC^k
\quad\text{for every }k\ge0.
}
\]

\begin{proof}
The base case holds by assumption. Suppose the bounds hold
for all indices up to \(k+1\). By the second hypothesis,
the recurrence relations imply the bounds at index \(k+2\).
Thus the bounds hold at every index by strong induction.
\end{proof}

\subsection{Geometric Majorant and Taylor Remainder}

Define
\[
M(r)=\sum_{k=0}^{\infty}KC^k r^k
=\frac{K}{1-Cr},\qquad |r|<\frac1C.
\]
The coefficient bounds imply absolute convergence of all
three power series for \(|r|<1/C\).

Let
\[
P_{a,10}(r)=\sum_{k=0}^{10}a_k r^k
\]
denote the tenth-order Taylor polynomial of \(a\), and
define \(P_{b,10}\) and \(P_{f,10}\) analogously. For
\(0\le\epsilon<1/C\),
\begin{align*}
|a(\epsilon)-P_{a,10}(\epsilon)|
&\le \sum_{k=11}^{\infty}|a_k|\epsilon^k\\
&\le K\sum_{k=11}^{\infty}(C\epsilon)^k\\
&=\frac{K(C\epsilon)^{11}}{1-C\epsilon}.
\end{align*}
The same argument applies to \(b\) and \(f\). Consequently,
\[
\boxed{
\begin{aligned}
|R_{a,10}(\epsilon)|&\le
\frac{K(C\epsilon)^{11}}{1-C\epsilon},\\
|R_{b,10}(\epsilon)|&\le
\frac{K(C\epsilon)^{11}}{1-C\epsilon},\\
|R_{f,10}(\epsilon)|&\le
\frac{K(C\epsilon)^{11}}{1-C\epsilon}.
\end{aligned}
}
\]
These estimates are rigorous provided the majorant hypotheses
have been verified for the actual recurrence and initial data.

% Inductive step needs more explanation: 

\end{document}